\rtmtight
\begin{tblr}{width=\textwidth,colspec={|Q[c,m,2.1cm]|X[l,m]|},hlines,vlines,hspan=minimal,rowsep=1.5pt,colsep=3pt,row{1}={bg=headgray,font=\bfseries}}
\SetCell{c,m}\hfil\logo\hfil & \textbf{POLITEKNIK NEGERI MALANG}\par\textbf{JURUSAN TEKNOLOGI INFORMASI}\par\textbf{PROGRAM STUDI: D4 TEKNIK INFORMATIKA} \\
\SetCell[c=2]{l} \textbf{RENCANA TUGAS MAHASISWA (RTM) DAN RUBRIK PENILAIAN} & \\
Nama Mata Kuliah & Teknologi Terapan \\
Kode Mata Kuliah & RTI256207 \quad \textbf{sks / Jam perminggu:} 6 SKS/12 jam \quad \textbf{Semester:} 6 \\
Dosen Pengampu & Tim Pengajar Program Studi D4 TI \\
\end{tblr}
\assessmentblock{Analisis Kebutuhan Teknologi Industri}{Case Method (CM) -- analisis dan perancangan terstruktur}{SCPMK0204-05601, SCPMK0801-05603}{Mahasiswa melakukan analisis kebutuhan teknologi industri secara komprehensif, mencakup needs assessment, technology mapping, feasibility study, dan penyusunan rencana implementasi teknologi.}{Melakukan observasi dan wawancara kebutuhan industri, memetakan teknologi yang tersedia, mengevaluasi kelayakan, menyusun rencana implementasi, dan mempresentasikan kepada stakeholder.}{Laporan analisis kebutuhan, technology canvas, feasibility study, rencana implementasi (roadmap dan manajemen risiko), dan presentasi kepada stakeholder.}{Kelengkapan needs assessment dan technology mapping & 6 & Kebutuhan industri diidentifikasi melalui metode yang tepat, teknologi yang relevan dipetakan, dan gap antara kondisi saat ini dengan target terdefinisi jelas. & Needs assessment dilakukan tetapi technology mapping atau gap analysis masih belum komprehensif. & Analisis kebutuhan tidak mendalam, technology mapping tidak dilakukan, atau gap tidak teridentifikasi. \\ \\
Kualitas feasibility study dan justifikasi pemilihan teknologi & 5 & Kelayakan teknis, finansial, dan operasional dievaluasi secara objektif; pemilihan teknologi didukung analisis trade-off yang kuat. & Feasibility study ada tetapi beberapa aspek kelayakan atau justifikasi pemilihan teknologi masih perlu penguatan. & Feasibility study tidak komprehensif atau pemilihan teknologi tidak didukung analisis yang memadai. \\ \\
Perencanaan implementasi dan manajemen risiko & 4 & Roadmap implementasi realistis, risiko teridentifikasi dan dimitigasi, dan rencana dikomunikasikan kepada stakeholder. & Perencanaan ada tetapi identifikasi risiko atau komunikasi kepada stakeholder masih kurang lengkap. & Tidak ada perencanaan yang terstruktur, risiko tidak diidentifikasi, atau rencana tidak dikomunikasikan. \\}

\assessmentblock{Sprint Implementasi Teknologi Terapan}{Project Based Learning (PBL) -- Sprint 1-2, UTS Milestone, Sprint 3-4}{SCPMK0503-05602, SCPMK0204-05601, SCPMK0801-05603}{Mahasiswa mengimplementasikan solusi teknologi terapan secara iteratif melalui sprint, mulai dari prototipe awal hingga integrasi dan optimasi, dengan melibatkan stakeholder di setiap tahapan.}{Melaksanakan sprint planning, membangun prototipe, memvalidasi dengan pengguna, mengintegrasikan dengan sistem yang ada, mengoptimasi, dan mendokumentasikan setiap sprint.}{Prototipe per sprint, kode sumber atau konfigurasi sistem, hasil validasi pengguna, laporan sprint, dan evidence UTS milestone.}{Kualitas prototipe dan iterasi pengembangan & 18 & Prototipe di setiap sprint menunjukkan kemajuan yang terukur, fitur sesuai kebutuhan industri, dan feedback pengguna diinkorporasikan secara konsisten. & Prototipe berkembang tetapi inkorporasi feedback atau kesesuaian dengan kebutuhan industri masih belum konsisten. & Prototipe tidak berkembang secara bermakna, feedback pengguna tidak diinkorporasikan, atau fitur tidak sesuai kebutuhan. \\ \\
Implementasi solusi dan integrasi sistem & 20 & Solusi diimplementasikan dengan kualitas kode/konfigurasi yang baik, integrasi dengan sistem yang ada berjalan lancar, dan tidak ada regresi fungsi. & Implementasi dan integrasi berjalan untuk sebagian besar komponen, tetapi beberapa masalah integrasi atau regresi belum sepenuhnya teratasi. & Implementasi tidak berhasil, integrasi menimbulkan banyak masalah, atau solusi tidak dapat diintegrasikan dengan sistem yang ada. \\ \\
Stakeholder engagement dan feedback incorporation & 17 & Stakeholder dilibatkan secara aktif di setiap sprint review, feedback terdokumentasi, diprioritaskan, dan tercermin dalam iterasi berikutnya. & Stakeholder dilibatkan tetapi feedback tidak selalu terdokumentasi atau diinkorporasikan dengan baik. & Stakeholder tidak dilibatkan secara bermakna atau feedback tidak diinkorporasikan dalam pengembangan. \\}

\assessmentblock{Laporan Implementasi Teknologi}{Case Method (CM) -- laporan teknis implementasi}{SCPMK0503-05602, SCPMK0204-05601}{Mahasiswa menyusun laporan implementasi teknologi yang komprehensif, mencakup deskripsi solusi, proses implementasi, hasil evaluasi, kendala yang dihadapi, dan rekomendasi pengembangan lanjutan.}{Mendokumentasikan seluruh proses implementasi dari analisis hingga deployment, mengevaluasi hasil terhadap target awal, dan menyusun rekomendasi berbasis pengalaman lapangan.}{Laporan implementasi lengkap, dokumentasi teknis (arsitektur, konfigurasi, manual), evaluasi hasil, dan rekomendasi pengembangan.}{Kelengkapan dokumentasi teknis implementasi & 4 & Dokumentasi mencakup arsitektur solusi, konfigurasi, panduan instalasi, dan manual pengguna yang dapat digunakan untuk pemeliharaan dan pengembangan lanjutan. & Dokumentasi teknis ada tetapi beberapa komponen seperti panduan instalasi atau manual pengguna masih belum lengkap. & Dokumentasi tidak lengkap sehingga tidak dapat digunakan untuk pemeliharaan atau pengembangan lanjutan. \\ \\
Evaluasi hasil versus kebutuhan awal & 3 & Hasil implementasi dievaluasi secara objektif terhadap kebutuhan awal, gap diidentifikasi, dan penyebab perbedaan dijelaskan. & Evaluasi dilakukan tetapi perbandingan dengan kebutuhan awal atau analisis gap masih kurang mendalam. & Tidak ada evaluasi yang bermakna atau hasil tidak dibandingkan dengan kebutuhan awal. \\ \\
Transfer knowledge dan manual penggunaan & 3 & Dokumentasi transfer knowledge lengkap, sesi pelatihan dilaksanakan, dan pengguna akhir mampu mengoperasikan solusi secara mandiri. & Transfer knowledge dilakukan tetapi dokumentasi atau efektivitas pelatihan masih belum optimal. & Tidak ada proses transfer knowledge atau pengguna akhir tidak dapat mengoperasikan solusi secara mandiri. \\}

\assessmentblock{UAS Presentasi Teknologi Terapan}{UAS presentasi dan demo kepada industri}{SCPMK0801-05603, SCPMK0503-05602}{Mahasiswa mempresentasikan solusi teknologi terapan yang diimplementasikan kepada stakeholder industri, mendemonstrasikan produk akhir, dan mempertahankan keputusan desain serta strategi implementasi yang dipilih.}{Mempersiapkan materi presentasi, mendemonstrasikan solusi kepada stakeholder, menjelaskan dampak dan nilai bisnis, serta menjawab pertanyaan mengenai keputusan implementasi.}{Slide presentasi, demo solusi yang berjalan, laporan akhir implementasi, dan dokumentasi dampak teknologi.}{Kualitas demo produk dan bukti implementasi & 8 & Solusi berjalan stabil dalam kondisi produksi, semua fitur utama terdemo dengan baik, dan dampak terhadap proses bisnis industri dapat diukur. & Solusi berjalan untuk sebagian besar fitur, tetapi beberapa aspek implementasi atau pengukuran dampak masih terbatas. & Solusi tidak berjalan dengan baik dalam kondisi produksi atau dampak implementasi tidak dapat dibuktikan. \\ \\
Kemampuan presentasi dan komunikasi kepada stakeholder & 7 & Presentasi terstruktur, menggunakan bahasa yang sesuai audiens industri, dan mampu menjelaskan nilai bisnis serta teknis dari solusi yang diimplementasikan. & Presentasi cukup baik tetapi komunikasi nilai bisnis atau penyesuaian bahasa untuk audiens industri masih perlu peningkatan. & Presentasi tidak terstruktur, sulit dipahami, atau tidak mampu mengkomunikasikan nilai dari solusi yang dikembangkan. \\ \\
Dampak solusi dan rekomendasi pengembangan & 5 & Dampak solusi terhadap efisiensi/efektivitas industri terdokumentasi, dan rekomendasi pengembangan lanjutan berbasis data implementasi disajikan dengan jelas. & Dampak solusi disebutkan tetapi belum terdokumentasi secara terukur atau rekomendasi pengembangan masih bersifat umum. & Tidak ada evaluasi dampak yang bermakna atau tidak ada rekomendasi pengembangan yang relevan. \\}

\begin{tblr}{width=\textwidth,colspec={|Q[l,m,3.2cm]|X[l,m]|},hlines,vlines,hspan=minimal,row{1-3}={bg=headgray},rowsep=1.5pt,colsep=3pt}
Jadwal Pelaksanaan: & Minggu 1--17 sesuai RPS. Setiap evaluasi memiliki RTM dan rubrik multi-kriteria. \\
Lain-Lain Yang Diperlukan: & LMS, repositori kode sumber, cloud platform, prototyping tools, perangkat lunak sesuai mata kuliah, dan perangkat presentasi. \\
Pustaka: & Mengacu pustaka utama dan pendukung pada bagian RPS. \\
\end{tblr}
